% Results of sec:method-validity (methodology/10b_feature_validity.tex). Every [..] is a value to fill in
% from the aggregate outputs on dtu-gpu-7 (counts 1-4 shown as <5); nothing here is to be guessed.
% Order follows the methods: technical accuracy, construct validity, sensitivity, variance
% explained by risk factors, decision. Repeatability is not assessed (discussion).
% Delete the template sentence that does not apply after each table is filled.
% Wide placeholder tables are scaled with \resizebox (graphicx); check them again once filled.
% Red clarification notes; needs \usepackage{xcolor} in the preamble. Remove before submission.
\providecommand{\clarify}[1]{\textcolor{red}{\textbf{Clarification needed:} #1}}
\section{Measurement Properties of the Geometric Features}
\label{sec:results-validity}

The measurement properties of Section~\ref{sec:method-validity} were assessed in [..] clean-baseline participants with a visual read, of whom [..] had a follow-up scan, and in [..] participants with plaque at baseline as a secondary analysis.

\subsection{Technical Accuracy on Synthetic Geometry}
\label{sec:results-validity-accuracy}

\begin{table}[htbp]
  \centering
  \caption{Technical accuracy of each definition on synthetic geometry with smooth centerline displacement and radius errors at an amplitude of [..]\,mm. Bias is relative to the definition's own value without errors; the scale-free error is the standard deviation (SD) of the error divided by the between-participant SD in the cohort. Primary definitions first.}
  \label{tab:results-validity-accuracy}
  \footnotesize
  \setlength{\tabcolsep}{3pt}
  \resizebox{\linewidth}{!}{%
  \begin{tabular}{llcccc}
    \toprule
    Feature & Definition & Bias & SD of error & Linearity slope & Scale-free error \\
    \midrule
    Curvature & mean per mm, 5\,mm chord & [..] & [..] & [..] & [..] \\
     & mean per mm, 2.5\,mm chord & [..] & [..] & [..] & [..] \\
     & mean per mm, 10\,mm chord & [..] & [..] & [..] & [..] \\
     & count of bends (sine curves) & [..] & [..] & [..] & [..] \\
     & 90th percentile & [..] & [..] & [..] & [..] \\
     & largest turning angle & [..] & [..] & [..] & [..] \\
    \addlinespace
    Taper & Theil--Sen, area-equivalent & [..] & [..] & [..] & [..] \\
     & end-to-end ratio & [..] & [..] & [..] & [..] \\
     & distance-transform radius & [..] & [..] & [..] & [..] \\
     & no exclusions & [..] & [..] & [..] & [..] \\
    \addlinespace
    Diameter ratios & area-equivalent, [..]\,mm & [..] & [..] & [..] & [..] \\
     & minimum diameter & [..] & [..] & [..] & [..] \\
    \addlinespace
    LM angle & fitted over [..]\,mm & [..] & [..] & [..] & [..] \\
     & fitted over 5\,mm & [..] & [..] & [..] & [..] \\
     & fitted over 10\,mm & [..] & [..] & [..] & [..] \\
    \bottomrule
  \end{tabular}
  }
\end{table}

The primary definitions had a scale-free error of [..] to [..] (Table~\ref{tab:results-validity-accuracy}).
Against the analytic curvature, the 5\,mm chord had a bias of [..] and resolved bends with a radius down to [..]\,mm.
With independent point noise as the worst case, the scale-free error rose to [..].
[One sentence for any alternative with a clearly lower scale-free error than the primary definition, and whether the primary was kept.]

\subsection{Construct Validity}
\label{sec:results-validity-groups}

\begin{table}[htbp]
  \centering
  \caption{Tests of the a priori hypotheses in the clean-baseline cohort. Estimates as standardized differences (SD) or degrees, adjusted for age, with 95\% confidence intervals (90\% for hypotheses of no difference) in brackets. Bands and evidence as in Table~\ref{tab:method-validity}.}
  \label{tab:results-validity-groups}
  \footnotesize
  \setlength{\tabcolsep}{2pt}
  \resizebox{\linewidth}{!}{%
  \begin{tabular}{lllllll}
    \toprule
    Feature & Hypothesis & Type & Band or margin & Confirmatory estimate & Secondary estimate & Confirmed \\
    \midrule
    Curvature & women vs men & K & 0.2 to 1.0 SD & [..] ([..], [..]) & [..] ([..], [..]) & [..] \\
     & hypertension & K & $>0$, below sex & [..] ([..], [..]) & [..] ([..], [..]) & [..] \\
     & LDL & D & $\pm 0.2$ SD & [..] ([..], [..]) & & [..] \\
     & smoking & D & $\pm 0.2$ SD & [..] ([..], [..]) & & [..] \\
     & family history & D & $\pm 0.2$ SD & [..] ([..], [..]) & & [..] \\
    \addlinespace
    Taper & diameter, men vs women & K & 5 to 30\% & [..] ([..], [..]) & [..] ([..], [..]) & [..] \\
     & fractional taper, sex & K & $\pm 0.2$ SD & [..] ([..], [..]) & [..] ([..], [..]) & [..] \\
     & fractional taper, BSA & K & $\pm 0.2$ SD & [..] ([..], [..]) & & [..] \\
     & LAD vs RCA & K & $\geq 0.2$ SD & [..] ([..], [..]) & & [..] \\
    \addlinespace
    LM angle & men vs women & K & 5 to 15$^{\circ}$ & [..] ([..], [..]) & [..] ([..], [..]) & [..] \\
     & ramus intermedius & K & 7 to 21$^{\circ}$ & [..] ([..], [..]) & & [..] \\
    \addlinespace
    Diameter ratios & left vs right dominance & K & $\geq 0.2$ SD & [..] ([..], [..]) & & [..] \\
     & women vs men & K & $\pm 0.2$ SD & [..] ([..], [..]) & [..] ([..], [..]) & [..] \\
     & Finet ratio (median) & I & 0.60 to 0.72 & [..] ([..], [..]) & & [..] \\
    \bottomrule
  \end{tabular}
  }
\end{table}

Of the [..] hypotheses, [..] were confirmed, including [..] of the [..] graded strong (Table~\ref{tab:results-validity-groups}).
Curvature was [..] standard deviations higher in women than in men (95\% CI [..] to [..]), [inside / outside] the band of 0.2 to 1.0; adjusted for body surface area, the difference was [..] (95\% CI [..] to [..]).
The LM bifurcation angle was [..]$^{\circ}$ wider in men (95\% CI [..] to [..]) and [..]$^{\circ}$ wider with a ramus intermedius (95\% CI [..] to [..]).
[All hypotheses of no difference were confirmed / The hypothesis of no difference for [..] was not confirmed, with a 90\% CI of [..] to [..]].
\clarify{one sentence per unconfirmed hypothesis graded strong: which feature, and what was examined in its extraction.}

\subsection{Sensitivity to Definitional Choices}
\label{sec:results-validity-sensitivity}

\begin{table}[htbp]
  \centering
  \caption{Confirmation of the hypotheses under each definition. Primary definitions first.}
  \label{tab:results-validity-sensitivity}
  \footnotesize
  \setlength{\tabcolsep}{3pt}
  \resizebox{\linewidth}{!}{%
  \begin{tabular}{llcc}
    \toprule
    Feature & Definition & Hypotheses confirmed & Strong confirmed \\
    \midrule
    Curvature & mean per mm, 5\,mm chord & [..] of [..] & [..] of [..] \\
     & mean per mm, 2.5\,mm chord & [..] of [..] & [..] of [..] \\
     & mean per mm, 10\,mm chord & [..] of [..] & [..] of [..] \\
     & count of bends & [..] of [..] & [..] of [..] \\
     & 90th percentile & [..] of [..] & [..] of [..] \\
     & largest turning angle & [..] of [..] & [..] of [..] \\
    \addlinespace
    Taper & Theil--Sen, area-equivalent & [..] of [..] & [..] of [..] \\
     & end-to-end ratio & [..] of [..] & [..] of [..] \\
     & distance-transform radius & [..] of [..] & [..] of [..] \\
     & no exclusions & [..] of [..] & [..] of [..] \\
    \addlinespace
    Diameter ratios & area-equivalent, [..]\,mm & [..] of [..] & [..] of [..] \\
     & other distances & [..] of [..] & [..] of [..] \\
     & minimum diameter & [..] of [..] & [..] of [..] \\
    \addlinespace
    LM angle & fitted over [..]\,mm & [..] of [..] & [..] of [..] \\
     & fitted over 5 / 10\,mm & [..] of [..] & [..] of [..] \\
     & ramus excluded & [..] of [..] & [..] of [..] \\
    \bottomrule
  \end{tabular}
  }
\end{table}

The higher curvature in women was [found under all / found under [..] of the [..]] definitions of curvature (Table~\ref{tab:results-validity-sensitivity}), [including / but not] at the 10\,mm chord, at which noise contributes least.
[List any hypothesis that depended on the definition.]

\subsection{Variance Explained by Known Risk Factors}
\label{sec:results-validity-split}

\begin{table}[htbp]
  \centering
  \caption{Share of the variance of each feature explained by the covariates of the base model ($R^2$).}
  \label{tab:results-validity-split}
  \footnotesize
  \begin{tabular}{lc}
    \toprule
    Feature & $R^2$ \\
    \midrule
    Curvature & [..] \\
    Taper & [..] \\
    LM bifurcation angle & [..] \\
    LM$\to$LAD ratio & [..] \\
    LM$\to$LCx ratio & [..] \\
    \bottomrule
  \end{tabular}
\end{table}

The base-model covariates explained between [..] and [..] of the variance of the features (Table~\ref{tab:results-validity-split}), leaving [most / much / little] of it for the geometry.

\subsection{Decision}
\label{sec:results-validity-decision}

[..] of the five features met the prespecified rules of Section~\ref{sec:method-validity-reading} and were retained unchanged.
[If a feature failed:] [..] did not meet the rule on [..]; [what was examined, and which alternative replaced the primary definition, if any].
